\documentclass[a4paper,12pt]{article}

% --- Подключение пакетов ---
\usepackage[T2A]{fontenc}
\usepackage[utf8]{inputenc}
\usepackage[russian]{babel}

\usepackage{amsmath}
\usepackage{amssymb}
\usepackage{amsfonts}

\usepackage{graphicx}
\usepackage{geometry}
\usepackage{listings}
\usepackage{xcolor}
\usepackage{hyperref}
\usepackage{float}
\usepackage{booktabs} % Для красивых таблиц
\usepackage{multirow} % Для объединения строк в таблицах
\usepackage{array}

% Колонка фиксированной ширины с центрированием и запасом по бокам —
% устраняет визуальное "слипание" значений и процентов в таблицах.
\newcolumntype{Y}{>{\centering\arraybackslash}p{1.75cm}}

% --- Настройка полей ---
\geometry{left=2cm, right=1.5cm, top=2cm, bottom=2cm}

% --- Настройка ссылок ---
\hypersetup{
    colorlinks=true,
    linkcolor=black,
    filecolor=magenta,      
    urlcolor=blue,
    pdftitle={УИР 1 - Моделирование},
}

% --- Настройка стиля для вставки кода ---
\lstdefinestyle{mypython}{
    language=Python,
    basicstyle=\ttfamily\scriptsize,
    keywordstyle=\color{blue}\bfseries,
    stringstyle=\color{red},
    commentstyle=\color{green!60!black},
    numbers=left,
    numberstyle=\tiny\color{gray},
    stepnumber=1,
    numbersep=8pt,
    showspaces=false,
    showstringspaces=false,
    tabsize=4,
    frame=single,
    rulecolor=\color{black},
    breaklines=true,
    breakatwhitespace=false,
    captionpos=b,
    xleftmargin=15pt,
    xrightmargin=5pt,
    framexleftmargin=12pt,
    backgroundcolor=\color{gray!5},
    literate=%
        {а}{{\cyra}}1 {б}{{\cyrb}}1 {в}{{\cyrv}}1 {г}{{\cyrg}}1 {д}{{\cyrd}}1
        {е}{{\cyre}}1 {ё}{{\cyrie}}1 {ж}{{\cyrzh}}1 {з}{{\cyrz}}1 {и}{{\cyri}}1
        {й}{{\cyrishrt}}1 {к}{{\cyrk}}1 {л}{{\cyrl}}1 {м}{{\cyrm}}1 {н}{{\cyrn}}1
        {о}{{\cyro}}1 {п}{{\cyrp}}1 {р}{{\cyrr}}1 {с}{{\cyrs}}1 {т}{{\cyrt}}1
        {у}{{\cyru}}1 {ф}{{\cyrf}}1 {х}{{\cyrh}}1 {ц}{{\cyrc}}1 {ч}{{\cyrch}}1
        {ш}{{\cyrsh}}1 {щ}{{\cyrshch}}1 {ъ}{{\cyrhrdsn}}1 {ы}{{\cyrery}}1
        {ь}{{\cyrsftsn}}1 {э}{{\cyrerev}}1 {ю}{{\cyryu}}1 {я}{{\cyrya}}1
        {А}{{\CYRA}}1 {Б}{{\CYRB}}1 {В}{{\CYRV}}1 {Г}{{\CYRG}}1 {Д}{{\CYRD}}1
        {Е}{{\CYRE}}1 {Ё}{{\CYRIE}}1 {Ж}{{\CYRZH}}1 {З}{{\CYRZ}}1 {И}{{\CYRI}}1
        {Й}{{\CYRISHRT}}1 {К}{{\CYRK}}1 {Л}{{\CYRL}}1 {М}{{\CYRM}}1 {Н}{{\CYRN}}1
        {О}{{\CYRO}}1 {П}{{\CYRP}}1 {Р}{{\CYRR}}1 {С}{{\CYRS}}1 {Т}{{\CYRT}}1
        {У}{{\CYRU}}1 {Ф}{{\CYRF}}1 {Х}{{\CYRH}}1 {Ц}{{\CYRC}}1 {Ч}{{\CYRCH}}1
        {Ш}{{\CYRSH}}1 {Щ}{{\CYRSHCH}}1 {Ъ}{{\CYRHRDSN}}1 {Ы}{{\CYRERY}}1
        {Ь}{{\CYRSFTSN}}1 {Э}{{\CYREREV}}1 {Ю}{{\CYRYU}}1 {Я}{{\CYRYA}}1
}
\lstset{style=mypython}

\begin{document}

% ===========================================================================
% 1. ТИТУЛЬНЫЙ ЛИСТ
% ===========================================================================
\begin{titlepage}
    \centering
    {\fontsize{11}{13}\selectfont \textbf{Министерство науки и высшего образования Российской Федерации}} \\
    {\fontsize{7}{10}\selectfont ФЕДЕРАЛЬНОЕ ГОСУДАРСТВЕННОЕ АВТОНОМНОЕ ОБРАЗОВАТЕЛЬНОЕ УЧРЕЖДЕНИЕ ВЫСШЕГО ОБРАЗОВАНИЯ} \\
    \vspace{0.2cm} 
    {\fontsize{14}{17}\selectfont \textbf{«НАЦИОНАЛЬНЫЙ ИССЛЕДОВАТЕЛЬСКИЙ УНИВЕРСИТЕТ ИТМО»}}
    
    \vspace{1cm} 
    
    {\fontsize{12}{14}\selectfont \textbf{Группа P3311}}
    
    \vspace{1cm}
    
    {\fontsize{12}{14}\selectfont \textbf{Дисциплина: Моделирование}}
    
    \vspace*{\fill}
    
    {\fontsize{18}{22}\selectfont \textbf{УЧЕБНО-ИССЛЕДОВАТЕЛЬСКАЯ РАБОТА 1 (УИР 1)}} \\
    \vspace{0.5cm}
    {\fontsize{14}{17}\selectfont \textbf{по теме: «Обработка результатов измерений: статистический анализ числовой последовательности»}}
    
    \vspace{2cm}
    
    {\fontsize{12}{14}\selectfont \textbf{Вариант 207}}
    
    \vspace{2cm}
    
    \begin{flushright}
        \fontsize{12}{14}\selectfont
        Выполнил: \\
        Михайлов Петр Сергеевич \\
        \vspace{1em}
        Преподаватель: \\
        Тропченко Андрей Александрович
    \end{flushright}
    
    \vfill
    
    {\fontsize{12}{14}\selectfont
    Санкт-Петербург, 2026
    }
\end{titlepage}

\setcounter{page}{2}

\newpage

% ===========================================================================
% ОГЛАВЛЕНИЕ
% ===========================================================================
\tableofcontents

\newpage

% ===========================================================================
% ОСНОВНАЯ ЧАСТЬ
% ===========================================================================

\section*{Цель работы}
\addcontentsline{toc}{section}{Цель работы}
Изучение методов обработки и статистического анализа результатов измерений на примере заданной числовой последовательности путем оценки числовых моментов и выявления свойств последовательности на основе корреляционного анализа, а также аппроксимация закона распределения заданной последовательности по двум числовым моментам случайной величины.

Исходный код для генерации расчетов доступен на GitHub: \href{https://github.com/ВАША_ССЫЛКА}{\textbf{Ссылка на репозиторий}}.

% ===========================================================================
% ХОД РАБОТЫ
% ===========================================================================
\section*{Ход работы}
\addcontentsline{toc}{section}{Ход работы}

% ------------------ ЭТАП 1 ------------------
\subsection*{Этап 1. Форма №1. Оценки математического ожидания и дисперсии}
\addcontentsline{toc}{subsection}{Этап 1. Форма №1}
Оценки числовых моментов заданной числовой последовательности и доверительные интервалы сведены в таблицу. Значения в процентах показывают относительные отклонения полученных значений от эталонных (рассчитанных для выборки $N=300$).

\begin{table}[H]
\centering
\caption{Форма 1. Характеристики заданной ЧП (вариант 207)}
\label{tab:form1}
\begin{tabular}{|c|c|Y|Y|Y|Y|Y|Y|}
\hline
\multicolumn{2}{|c|}{\textbf{Характеристика}} & \multicolumn{6}{c|}{\textbf{Количество случайных величин}} \\ \cline{3-8} 
\multicolumn{2}{|c|}{} & \textbf{10} & \textbf{20} & \textbf{50} & \textbf{100} & \textbf{200} & \textbf{300} \\ \hline
\multirow{2}{*}{Мат.ож.} & Знач. & 9.800 & 7.650 & 10.660 & 14.560 & 11.200 & 11.517 \\ \cline{2-8} 
 & \% & 14.91 & 33.57 & 7.44 & 26.43 & 2.75 & 0.00 \\ \hline
\multirow{2}{*}{Дов. инт. (0,9)} & Знач. & $\pm$9.045 & $\pm$5.145 & $\pm$4.382 & $\pm$5.494 & $\pm$2.998 & $\pm$2.653 \\ \cline{2-8} 
 & \% & 240.88 & 93.91 & 65.15 & 107.06 & 12.98 & 0.00 \\ \hline
\multirow{2}{*}{Дов. инт. (0,95)} & Знач. & $\pm$10.790 & $\pm$6.138 & $\pm$5.227 & $\pm$6.554 & $\pm$3.576 & $\pm$3.165 \\ \cline{2-8} 
 & \% & 240.88 & 93.91 & 65.15 & 107.06 & 12.98 & 0.00 \\ \hline
\multirow{2}{*}{Дов. инт. (0,99)} & Знач. & $\pm$14.181 & $\pm$8.067 & $\pm$6.870 & $\pm$8.614 & $\pm$4.700 & $\pm$4.160 \\ \cline{2-8} 
 & \% & 240.88 & 93.91 & 65.15 & 107.06 & 12.98 & 0.00 \\ \hline
\multirow{2}{*}{Дисперсия} & Знач. & 303.067 & 196.134 & 355.658 & 1118.188 & 665.879 & 782.431 \\ \cline{2-8} 
 & \% & 61.27 & 74.93 & 54.54 & 42.91 & 14.90 & 0.00 \\ \hline
\multirow{2}{*}{С.к.о.} & Знач. & 17.409 & 14.005 & 18.859 & 33.439 & 25.805 & 27.972 \\ \cline{2-8} 
 & \% & 37.76 & 49.93 & 32.58 & 19.55 & 7.75 & 0.00 \\ \hline
\multirow{2}{*}{К-т вариации} & Знач. & 1.776 & 1.831 & 1.769 & 2.297 & 2.304 & 2.429 \\ \cline{2-8} 
 & \% & 26.86 & 24.63 & 27.16 & 5.44 & 5.14 & 0.00 \\ \hline
\end{tabular}
\end{table}

\textbf{Вывод из 1 этапа:} Дисперсия и среднеквадратическое отклонение имеют высокую относительную погрешность на малых объемах выборки, но при объеме от 200 значений наблюдается стабилизация параметров. Коэффициент вариации для эталонной выборки составляет $V = 2.429$, что превышает единицу и свидетельствует об очень высокой степени рассеивания данных относительно математического ожидания.


% ------------------ ЭТАП 2 ------------------
\subsection*{Этап 2. График №1. Значения заданной числовой последовательности}
\addcontentsline{toc}{subsection}{Этап 2. График №1}
\begin{figure}[H]
    \centering
    \includegraphics[width=0.75\textwidth]{graph_1.png}
    \caption{График 1. Значения заданной числовой последовательности}
\end{figure}

\textbf{Вывод из 2 этапа:} Изучив график, можно сделать вывод, что исходная последовательность не является ни периодической, ни монотонно возрастающей, ни убывающей. Присутствуют резкие хаотичные выбросы, что указывает на стохастическую природу данных.


% ------------------ ЭТАП 3 ------------------
\subsection*{Этап 3. Форма 3. Результаты автокорреляционного анализа}
\addcontentsline{toc}{subsection}{Этап 3. Форма 3}
Результаты представлены как в числовом (Таблица \ref{tab:form3_1}), так и в графическом виде.

\begin{table}[H]
\centering
\caption{Форма 3 (исходная ЧП). Коэффициенты автокорреляции}
\label{tab:form3_1}
\resizebox{\textwidth}{!}{
\begin{tabular}{|l|c|c|c|c|c|c|c|c|c|c|}
\hline
\textbf{Сдвиг ЧП} & \textbf{1} & \textbf{2} & \textbf{3} & \textbf{4} & \textbf{5} & \textbf{6} & \textbf{7} & \textbf{8} & \textbf{9} & \textbf{10} \\ \hline
\textbf{К-т АК} & 0.043 & -0.067 & -0.024 & 0.008 & -0.044 & 0.059 & 0.043 & -0.074 & 0.028 & 0.032 \\ \hline
\end{tabular}
}
\end{table}

\begin{figure}[H]
    \centering
    \includegraphics[width=0.7\textwidth]{graph_acf_orig.png}
    \caption{Автокорреляционный анализ исходной ЧП}
\end{figure}

\textbf{Вывод из 3 этапа:} Последовательность можно считать случайной, так как данные коэффициенты (и их график) указывают на то, что между значениями выборки нет линейной зависимости, тенденции и периодичности (все значения колеблются в узком коридоре около нуля).


% ------------------ ЭТАП 4 ------------------
\subsection*{Этап 4. График 2. Гистограмма распределения частот}
\addcontentsline{toc}{subsection}{Этап 4. График 2}
\begin{figure}[H]
    \centering
    \includegraphics[width=0.75\textwidth]{graph_2.png}
    \caption{График 2. Гистограмма распределения частот}
    \label{fig:graph2}
\end{figure}

Для наглядности хвоста распределения та же гистограмма приведена ниже с логарифмической шкалой по оси ординат (Рис. \ref{fig:graph2log}) — на ней хорошо видно, что немногочисленные крупные выбросы образуют плавно убывающий "хвост", а не единичные случайные всплески.

\begin{figure}[H]
    \centering
    \includegraphics[width=0.75\textwidth]{graph_2_log.png}
    \caption{График 2 (логарифмическая шкала по оси Y)}
    \label{fig:graph2log}
\end{figure}

Численные значения гистограммы (границы интервалов и частоты попадания) сведены в Таблицу \ref{tab:hist_orig}.

\begin{table}[H]
\centering
\caption{Группированный статистический ряд исходной ЧП (25 интервалов)}
\label{tab:hist_orig}
\resizebox{\textwidth}{!}{
\begin{tabular}{|c|c|c|c||c|c|c|c|}
\hline
\textbf{№} & \textbf{Лев. гр.} & \textbf{Прав. гр.} & \textbf{Частота} & \textbf{№} & \textbf{Лев. гр.} & \textbf{Прав. гр.} & \textbf{Частота} \\ \hline
1 & 0.000 & 11.200 & 247 & 14 & 145.600 & 156.800 & 0 \\ \hline
2 & 11.200 & 22.400 & 15 & 15 & 156.800 & 168.000 & 0 \\ \hline
3 & 22.400 & 33.600 & 10 & 16 & 168.000 & 179.200 & 0 \\ \hline
4 & 33.600 & 44.800 & 7 & 17 & 179.200 & 190.400 & 2 \\ \hline
5 & 44.800 & 56.000 & 3 & 18 & 190.400 & 201.600 & 0 \\ \hline
6 & 56.000 & 67.200 & 8 & 19 & 201.600 & 212.800 & 0 \\ \hline
7 & 67.200 & 78.400 & 2 & 20 & 212.800 & 224.000 & 0 \\ \hline
8 & 78.400 & 89.600 & 0 & 21 & 224.000 & 235.200 & 0 \\ \hline
9 & 89.600 & 100.800 & 3 & 22 & 235.200 & 246.400 & 0 \\ \hline
10 & 100.800 & 112.000 & 1 & 23 & 246.400 & 257.600 & 0 \\ \hline
11 & 112.000 & 123.200 & 1 & 24 & 257.600 & 268.800 & 0 \\ \hline
12 & 123.200 & 134.400 & 0 & 25 & 268.800 & 280.000 & 1 \\ \hline
13 & 134.400 & 145.600 & 0 & & & & \\ \hline
\end{tabular}
}
\end{table}


\textbf{Вывод из 4 этапа:} Исходя из гистограммы, можно видеть, что подавляющая часть значений располагается в промежутке вблизи нуля. На Графике 2 сам хвост распределения визуально слабо различим: единичные выбросы (вплоть до значений $250+$) растягивают ось $X$, из-за чего столбцы, соответствующие хвосту, "прижимаются" к нулевой линии по оси $Y$. Логарифмическая шкала на Рис. \ref{fig:graph2log} устраняет эту проблему и подтверждает наличие протяжённого убывающего хвоста. Совместно с ранее найденным коэффициентом вариации $V > 1$, можно строго предположить, что закон распределения ЧП — гиперэкспоненциальный.


% ------------------ ЭТАП 5 ------------------
\subsection*{Этап 5. Параметры аппроксимирующего закона распределения}
\addcontentsline{toc}{subsection}{Этап 5. Аппроксимация}
Для данной по варианту выборки коэффициент вариации больше единицы ($v = 2.429$). Следовательно, для аппроксимации последовательности будет использован гиперэкспоненциальный закон. Значения математического ожидания $t = M(X) = 11.517$ и коэффициента вариации были определены ранее.

Ограничение для параметра $q$:
$$ q \le \frac{2}{1 + v^2} \approx 0.2899 $$
Выберем $q$ равным 80\% от максимума для адекватного разброса: $q = 0.2319$. Тогда параметры $t_1$ и $t_2$ находятся как:
$$ t_1 = t \left[ 1 + \sqrt{\frac{1-q}{2q}(v^2 - 1)} \right] = 44.3201 $$
$$ t_2 = t \left[ 1 - \sqrt{\frac{q}{2(1-q)}(v^2 - 1)} \right] = 1.6122 $$

Соответственно, мы получаем аппроксимирующий закон как композицию двух экспоненциальных распределений с вероятностями $q$ и $1-q$.

\textbf{Вывод из 5 этапа:} На основе аналитических соотношений для первых двух моментов были найдены параметры гиперэкспоненциального закона ($q$, $t_1$, $t_2$), которые полностью определяют математическую модель для последующей генерации.


% ------------------ ЭТАП 6 ------------------
\subsection*{Этап 6. Описание алгоритма формирования закона распределения}
\addcontentsline{toc}{subsection}{Этап 6. Алгоритм генерации}
Для генерации случайной последовательности используется язык программирования Python. Алгоритм базируется на методе обратной функции.
\begin{enumerate}
    \item Генерируется случайное число $U_1 \sim \mathcal{U}(0,1)$.
    \item Производится выбор ветви: если $U_1 < q$, то используется параметр $\lambda = 1/t_1$. В противном случае используется $\lambda = 1/t_2$.
    \item Генерируется второе случайное число $U_2 \sim \mathcal{U}(0,1)$.
    \item Случайная величина $X$ рассчитывается как $X = -\frac{1}{\lambda} \ln(U_2)$.
\end{enumerate}

\textbf{Вывод из 6 этапа:} Был успешно описан и программно реализован алгоритм формирования случайной числовой последовательности (ЧП), распределенной по гиперэкспоненциальному закону на основе найденных параметров.


% ------------------ ЭТАП 7 ------------------
\subsection*{Этап 7. Сравнение сгенерированной ЧП с исходной}
\addcontentsline{toc}{subsection}{Этап 7. Сравнение ЧП}

\begin{table}[H]
\centering
\caption{Форма 2. Характеристики сгенерированной случайной ЧП}
\label{tab:form2}
\begin{tabular}{|c|c|Y|Y|Y|Y|Y|Y|}
\hline
\multicolumn{2}{|c|}{\textbf{Характеристика}} & \multicolumn{6}{c|}{\textbf{Количество случайных величин}} \\ \cline{3-8} 
\multicolumn{2}{|c|}{} & \textbf{10} & \textbf{20} & \textbf{50} & \textbf{100} & \textbf{200} & \textbf{300} \\ \hline
\multirow{2}{*}{Мат.ож.} & Знач. & 27.136 & 22.138 & 20.655 & 14.759 & 12.425 & 13.767 \\ \cline{2-8} 
 & \% & 176.90 & 189.38 & 93.76 & 1.37 & 10.94 & 19.54 \\ \hline
\multirow{2}{*}{Дов. инт. (0,9)} & Знач. & $\pm$27.308 & $\pm$17.031 & $\pm$10.109 & $\pm$5.615 & $\pm$3.423 & $\pm$2.998 \\ \cline{2-8} 
 & \% & 201.91 & 231.02 & 130.70 & 2.21 & 14.18 & 13.00 \\ \hline
\multirow{2}{*}{Дов. инт. (0,95)} & Знач. & $\pm$32.576 & $\pm$20.317 & $\pm$12.060 & $\pm$6.699 & $\pm$4.084 & $\pm$3.577 \\ \cline{2-8} 
 & \% & 201.91 & 231.02 & 130.70 & 2.21 & 14.18 & 13.00 \\ \hline
\multirow{2}{*}{Дов. инт. (0,99)} & Знач. & $\pm$42.815 & $\pm$26.703 & $\pm$15.850 & $\pm$8.804 & $\pm$5.367 & $\pm$4.701 \\ \cline{2-8} 
 & \% & 201.91 & 231.02 & 130.70 & 2.21 & 14.18 & 13.00 \\ \hline
\multirow{2}{*}{Дисперсия} & Знач. & 2762.426 & 2149.076 & 1892.863 & 1168.051 & 868.159 & 999.132 \\ \cline{2-8} 
 & \% & 811.49 & 995.72 & 432.21 & 4.46 & 30.38 & 27.70 \\ \hline
\multirow{2}{*}{С.к.о.} & Знач. & 52.559 & 46.358 & 43.507 & 34.177 & 29.465 & 31.609 \\ \cline{2-8} 
 & \% & 201.91 & 231.02 & 130.70 & 2.21 & 14.18 & 13.00 \\ \hline
\multirow{2}{*}{К-т вариации} & Знач. & 1.937 & 2.094 & 2.106 & 2.316 & 2.371 & 2.296 \\ \cline{2-8} 
 & \% & 9.03 & 14.39 & 19.06 & 0.83 & 2.93 & 5.47 \\ \hline
\end{tabular}
\end{table}

\begin{figure}[H]
    \centering
    
    \includegraphics[width=0.8\textwidth]{graph_acf_gen.png}
    \caption{Автокорреляция сгенерированной ЧП}

\end{figure}

Численные значения коэффициентов автокорреляции сгенерированной ЧП, а также их сравнение с коэффициентами исходной ЧП, сведены в таблицу по форме 3.

\begin{table}[H]
\centering
\caption{Форма 3 (сравнение). Коэффициенты автокорреляции исходной и сгенерированной ЧП}
\label{tab:form3_2}
\resizebox{\textwidth}{!}{
\begin{tabular}{|l|c|c|c|c|c|c|c|c|c|c|}
\hline
\textbf{Сдвиг ЧП} & \textbf{1} & \textbf{2} & \textbf{3} & \textbf{4} & \textbf{5} & \textbf{6} & \textbf{7} & \textbf{8} & \textbf{9} & \textbf{10} \\ \hline
\textbf{К-т АК задан.} & 0.043 & -0.067 & -0.024 & 0.008 & -0.044 & 0.059 & 0.043 & -0.074 & 0.028 & 0.032 \\ \hline
\textbf{К-т АК сгенерир.} & 0.066 & 0.110 & -0.023 & -0.073 & 0.028 & -0.022 & -0.038 & 0.013 & -0.025 & -0.054 \\ \hline
\textbf{\%} & 53.49 & 264.03 & 5.83 & 1017.50 & 163.41 & 137.80 & 189.07 & 117.30 & 187.86 & 270.00 \\ \hline
\end{tabular}
}
\end{table}

Как видно из таблицы, относительные отклонения (\%) коэффициентов автокорреляции сгенерированной ЧП от исходной формально очень велики (десятки и сотни процентов). Это не свидетельствует о плохом качестве генерации: поскольку сами коэффициенты автокорреляции обеих последовательностей близки к нулю, любое малое абсолютное отклонение (порядка $0.02$--$0.05$) даёт большую относительную погрешность. Содержательный вывод здесь делается не по проценту, а по факту того, что оба ряда коэффициентов лежат в узком коридоре около нуля без выраженного тренда — то есть обе последовательности статистически независимы (некоррелированы) внутри себя.
    
\vspace{0.5cm}

\begin{figure}[H]
    \centering
    
    \includegraphics[width=0.8\textwidth]{graph_3.png}
    \caption{Гистограмма сгенерированной ЧП}
    
\end{figure}

Обе гистограммы (График 2 и График 3) построены в одинаковом нормированном масштабе — по оси ординат отложена плотность вероятности (\texttt{density=True}), что делает их визуальное сравнение корректным. Численные значения гистограммы сгенерированной ЧП (границы интервалов и частоты попадания) сведены в Таблицу \ref{tab:hist_gen}.

\begin{table}[H]
\centering
\caption{Группированный статистический ряд сгенерированной ЧП (25 интервалов)}
\label{tab:hist_gen}
\resizebox{\textwidth}{!}{
\begin{tabular}{|c|c|c|c||c|c|c|c|}
\hline
\textbf{№} & \textbf{Лев. гр.} & \textbf{Прав. гр.} & \textbf{Частота} & \textbf{№} & \textbf{Лев. гр.} & \textbf{Прав. гр.} & \textbf{Частота} \\ \hline
1 & 0.015 & 7.698 & 237 & 14 & 99.894 & 107.577 & 4 \\ \hline
2 & 7.698 & 15.381 & 8 & 15 & 107.577 & 115.260 & 0 \\ \hline
3 & 15.381 & 23.064 & 5 & 16 & 115.260 & 122.943 & 2 \\ \hline
4 & 23.064 & 30.747 & 7 & 17 & 122.943 & 130.626 & 0 \\ \hline
5 & 30.747 & 38.430 & 8 & 18 & 130.626 & 138.309 & 1 \\ \hline
6 & 38.430 & 46.113 & 5 & 19 & 138.309 & 145.992 & 1 \\ \hline
7 & 46.113 & 53.796 & 7 & 20 & 145.992 & 153.675 & 1 \\ \hline
8 & 53.796 & 61.479 & 2 & 21 & 153.675 & 161.358 & 2 \\ \hline
9 & 61.479 & 69.162 & 1 & 22 & 161.358 & 169.041 & 0 \\ \hline
10 & 69.162 & 76.845 & 3 & 23 & 169.041 & 176.724 & 0 \\ \hline
11 & 76.845 & 84.528 & 1 & 24 & 176.724 & 184.407 & 0 \\ \hline
12 & 84.528 & 92.211 & 2 & 25 & 184.407 & 192.090 & 2 \\ \hline
13 & 92.211 & 99.894 & 1 & & & & \\ \hline
\end{tabular}
}
\end{table}

Коэффициент корреляции между двумя числовыми последовательностями рассчитывался по формуле Пирсона:
$$ r = \frac{\sum (x_i - \bar{x})(y_i - \bar{y})}{\sqrt{\sum (x_i - \bar{x})^2 \times \sum (y_i - \bar{y})^2}} = -0.0510 $$

Для сгенерированной и полученной последовательности мы рассчитали корреляционную зависимость. Как мы видим из $|r| \approx 0.05$, корреляции между исходной и сгенерированной случайной последовательностями нет.

Отдельно стоит прокомментировать аномально большие относительные отклонения дисперсии в Таблице \ref{tab:form2} при малых объёмах выборки ($995.72\%$ при $N=20$ и $811.49\%$ при $N=10$). Это объясняется не ошибкой расчёта, а самой природой гиперэкспоненциального закона: он обладает ``тяжёлым хвостом'' -- подавляющее большинство наблюдений малы (генерируются с параметром $t_2 = 1.61$), но с относительно небольшой вероятностью $q = 0.23$ появляются значения из ``медленной'' ветви с $t_1 = 44.32$, которые вносят основной вклад именно во второй момент (дисперсию), а не в среднее. При выборке из $10$--$20$ значений скорее всего либо ни одного, либо всего одно такое ``хвостовое'' наблюдение не попадает в выборку (или попадает избыточно), из-за чего оценка дисперсии по малой выборке крайне неустойчива и может отличаться от эталонной в разы. Это же объясняет, почему для среднего (мат. ожидания) относительное отклонение при малых $N$ на порядок меньше дисперсии: первый момент гиперэкспоненциального распределения гораздо устойчивее к редким выбросам, чем второй.

\textbf{Вывод из 7 этапа:} При визуальном сравнении полученной гистограммы сгенерированной ЧП (График 3) с гистограммой исходной ЧП (График 2) видно, что их формы совпадают: большая часть значений сконцентрирована в начале, присутствует характерный длинный хвост. Тем самым доказано, что выбранная нами аппроксимация гиперэкспоненциальным законом подходит. Сравнение числовых характеристик исходной и сгенерированной ЧП (в частности, коэффициента вариации, а также значений при $N \geq 100$) также показало их достаточное сходство, а коэффициенты автокорреляции обеих последовательностей подтвердили отсутствие внутренней линейной зависимости.


% ------------------ ЭТАП 8 ------------------
\subsection*{8. Выводы}
\addcontentsline{toc}{subsection}{Выводы}
В рамках лабораторной работы была дана числовая последовательность, для которой мы определили математическое ожидание, дисперсию и другие параметры. Далее мы проанализировали построенную гистограмму, по которой не было выявлено возрастания, убывания или периодичности последовательности. Исследуемую последовательность можно назвать случайной исходя из автокорреляционного анализа. 

Затем мы вычислили параметры аппроксимирующего закона (гиперэкспоненциального, так как $V > 1$) и по ним сгенерировали новую последовательность. Коэффициенты вариации первой и второй последовательности близки ($2.42$ и $2.29$). Коэффициенты автокорреляции первой и второй последовательности варьируются около нуля, исходя из этого можно сделать вывод о том, что обе выборки независимы и случайны. Математическое ожидание и дисперсия отличаются, но это отличие лежит в рамках статистической погрешности для выборок конечного объема. Таким образом, подобранная модель адекватно описывает выданную последовательность.

\end{document}